%% ──────────────────────────────────────────────────────────────
%% templates_latex/resume.tex — Jinja2 template that renders to LaTeX
%% ──────────────────────────────────────────────────────────────
%%
%% CUSTOM JINJA2 DELIMITERS
%% ────────────────────────
%% This template does NOT use standard {{ }}/{% %} Jinja2 syntax.
%% Instead it uses delimiters that look like LaTeX commands:
%%
%%   \VAR{ variable }       — insert a variable's value
%%   \BLOCK{ if cond }      — start a logic block
%%   \BLOCK{ endif }        — end a logic block
%%   \BLOCK{ for x in y }   — loop
%%   \BLOCK{ endfor }       — end loop
%%   \COMMENT{ note }       — template comment (stripped from output)
%%
%% WHY?  LaTeX uses { } everywhere (\textbf{bold}, \begin{itemize}).
%% If we used Jinja2's default {{ }}, the template engine would
%% misparse LaTeX braces as template tags, causing cryptic errors.
%% The \COMMAND{...} convention is unambiguous because no real LaTeX
%% command is named \VAR, \BLOCK, or \COMMENT.
%%
%% ALL USER DATA IS PRE-ESCAPED
%% ────────────────────────────
%% Before this template is rendered, every string from the user's
%% profile has already been run through escape_latex(), which
%% replaces LaTeX special characters (& % $ # _ { } ~ ^ \) with
%% safe equivalents (\& \% \$ etc.).  This means we can safely
%% inject \VAR{ name } without worrying about compilation errors
%% from characters like & or #.
%% ──────────────────────────────────────────────────────────────

\documentclass[11pt, letterpaper]{article}

%% ── Packages ─────────────────────────────────────────────────
\usepackage[utf8]{inputenc}          % UTF-8 input encoding
\usepackage[T1]{fontenc}             % proper output glyph encoding
\usepackage[scaled=0.92]{helvet}     % Helvetica font (clean sans-serif)
\renewcommand{\familydefault}{\sfdefault} % Set sans-serif as default
\usepackage[margin=0.5in]{geometry}  % tight but readable margins
\usepackage{enumitem}                % customise list spacing
\usepackage{titlesec}                % customise section headings
\usepackage{hyperref}                % clickable links
\usepackage{xcolor}                  % colours for links
\usepackage{marvosym}            % Classic icons for contact info
\usepackage[normalem]{ulem}          % For clean underlines

%% ── Hyperlink styling ────────────────────────────────────────
\hypersetup{
    colorlinks=true,
    linkcolor=black,
    urlcolor=blue!70!black,
    pdfborder={0 0 0},
}

%% ── Section heading format ───────────────────────────────────
%% \titleformat{command}{format}{label}{sep}{before-code}[after-code]
%% We use a simple bold heading (without small caps) with a rule underneath.
\titleformat{\section}
  {\large\bfseries\raggedright}          % format
  {}                                      % label (none)
  {0em}                                   % separator
  {}                                      % before-code
  [\titlerule]                            % after-code: horizontal rule

\BLOCK{ if compact }%
\titlespacing*{\section}{0pt}{6pt}{2pt}%
\setlist[itemize]{label=\textbullet, nosep, leftmargin=1.5em, topsep=0pt}%
\newcommand{\entrysep}{\vspace{2pt}}%
\BLOCK{ else }%
\titlespacing*{\section}{0pt}{10pt}{4pt}%
\setlist[itemize]{label=\textbullet, nosep, leftmargin=1.5em, topsep=2pt}%
\newcommand{\entrysep}{\vspace{4pt}}%
\BLOCK{ endif }%

%% ── Remove page numbers ─────────────────────────────────────
\pagestyle{empty}

%% ── Zero Paragraph Indentation ──────────────────────────────
\setlength{\parindent}{0pt}

%% ── Unicode Rupee Symbol Mapping ────────────────────────────
\usepackage{tfrupee}
\usepackage{newunicodechar}
\newunicodechar{₹}{\rupee}

%% ══════════════════════════════════════════════════════════════
%%  DOCUMENT BEGIN
%% ══════════════════════════════════════════════════════════════
\begin{document}

%% ── Header ───────────────────────────────────────────────────
\begin{center}
    {\LARGE\bfseries \VAR{ personal_info.name }}\\[5pt]
    \small
    \BLOCK{ if personal_info.phone }%
        \Mobilefone\ \VAR{ personal_info.phone }%
    \BLOCK{ endif }%
    \BLOCK{ if personal_info.email }%
        \BLOCK{ if personal_info.phone } \quad \BLOCK{ endif }%
        \Letter\ \href{mailto:\VAR{ personal_info.email }}{\VAR{ personal_info.email }}%
    \BLOCK{ endif }%
    \BLOCK{ for label, url in personal_info.links.items() }%
        \quad \textbullet\ \href{\VAR{ url }}{\VAR{ label }}%
    \BLOCK{ endfor }
\end{center}

\vspace{-4pt}


\BLOCK{ if education }%
%% ── Education ────────────────────────────────────────────────
\section{Education}
\BLOCK{ for edu in education }%
\textbf{\VAR{ edu.institution }} \hfill \VAR{ edu.start_date } -- \VAR{ edu.end_date } \\
\textit{\VAR{ edu.degree }}\BLOCK{ if edu.gpa } \textit{(CGPA: \VAR{ edu.gpa })}\BLOCK{ endif }%
\BLOCK{ if edu.coursework }%
 \hfill \VAR{ edu.coursework | join(', ') }%
\BLOCK{ endif }%
\BLOCK{ if not loop.last }

\entrysep
\BLOCK{ endif }%
\BLOCK{ endfor }%
\BLOCK{ endif }


\BLOCK{ if experience }%
%% ── Experience ───────────────────────────────────────────────
\section{Experience}
\BLOCK{ for exp in experience }%
\textbf{\VAR{ exp.company }} \hfill \VAR{ exp.start_date } -- \VAR{ exp.end_date } \\
\textit{\VAR{ exp.role }}\BLOCK{ if exp.technologies } \ | \ \textit{\VAR{ exp.technologies | join(', ') }}\BLOCK{ endif }%
\BLOCK{ if exp.work_mode }%
 \hfill \VAR{ exp.work_mode }%
\BLOCK{ endif }%
\BLOCK{ if exp.bullets }
\begin{\VAR{ exp.bullets_type }}\BLOCK{ if exp.bullets_type == 'itemize' }[label=\textasteriskcentered, leftmargin=1.5em]\BLOCK{ else }[leftmargin=1.5em]\BLOCK{ endif }
\BLOCK{ for bullet in exp.bullets }
    \item \VAR{ bullet }
\BLOCK{ endfor }
\end{\VAR{ exp.bullets_type }}
\BLOCK{ endif }%
\BLOCK{ if not loop.last }
\entrysep
\BLOCK{ endif }%
\BLOCK{ endfor }%
\BLOCK{ endif }


\BLOCK{ if projects }%
%% ── Projects ─────────────────────────────────────────────────
\section{Projects}
\BLOCK{ for proj in projects }%
\textbf{\VAR{ proj.name }}\BLOCK{ if proj.link } \ | \ \href{\VAR{ proj.link }}{\textbf{\BLOCK{ if 'github' in proj.link.lower() }(GitHub)\BLOCK{ else }(Link)\BLOCK{ endif }}}\BLOCK{ endif }\BLOCK{ if proj.technologies } \ | \ \textit{\VAR{ proj.technologies | join(', ') }}\BLOCK{ endif }%
\BLOCK{ if proj.description and proj.description.strip() }%
\\ \VAR{ proj.description }%
\BLOCK{ endif }%
\BLOCK{ if proj.bullets }
\begin{\VAR{ proj.bullets_type }}\BLOCK{ if proj.bullets_type == 'itemize' }[label=\textasteriskcentered, leftmargin=1.5em]\BLOCK{ else }[leftmargin=1.5em]\BLOCK{ endif }
\BLOCK{ for bullet in proj.bullets }
    \item \VAR{ bullet }
\BLOCK{ endfor }
\end{\VAR{ proj.bullets_type }}
\BLOCK{ endif }%
\BLOCK{ if not loop.last }
\entrysep
\BLOCK{ endif }%
\BLOCK{ endfor }%
\BLOCK{ endif }


\BLOCK{ if skills and skills.categories }%
%% ── Skills ───────────────────────────────────────────────────
\section{Technical Skills}
\BLOCK{ for category, skill_list in skills.categories.items() }%
\textbf{\VAR{ category }:} \VAR{ skill_list | join(', ') }\BLOCK{ if not loop.last } \\ \BLOCK{ endif }%
\BLOCK{ endfor }%
\BLOCK{ endif }


\BLOCK{ if certifications }%
%% ── Certifications ───────────────────────────────────────────
\section{Certifications}
\BLOCK{ for cert in certifications }%
\textbf{\VAR{ cert.name }} — \VAR{ cert.issuer } \hfill \VAR{ cert.date }%
\BLOCK{ if cert.description }%
\\ \small{\VAR{ cert.description }}%
\BLOCK{ endif }%
\BLOCK{ if not loop.last }

\BLOCK{ endif }%
\BLOCK{ endfor }%
\BLOCK{ endif }


\BLOCK{ if achievements }%
%% ── Achievements ─────────────────────────────────────────────
\section{Achievements}
\VAR{ achievements }
\BLOCK{ endif }

\end{document}
